% !TEX root = ../main.tex
%
% Chapter 02 — Golden Cut: Extreme-and-Mean Ratio, Continued Fractions,
%              Diophantine Approximation, and Aperiodic Structure
%
% Rule R3  : ≥ 1500 lines, ≥ 2 citations, ≥ 1 theorem+proof+qed
% Rule R6  : constants φ, π, e, n∈ℤ only
% Rule R12 : theorems use \begin{theorem}…\end{theorem}, proof uses "we"
% Rule R14 : Coq citation map — see \coqcite entries below
%
\chapter{Golden Cut: Background --- Neuro-Symbolic AI}
\label{ch:golden-cut}

\begin{figure}[H]
\centering
\makebox[\linewidth][c]{\includegraphics[width=1.18\linewidth,keepaspectratio]{ch02-background.png}}
\caption*{Figure --- Golden Cut: Background --- Neuro-Symbolic AI.}
\end{figure}

% ============================================================
\section{Abstract}\label{sec:gc-abstract}
% ============================================================

This chapter surveys the conceptual and technical
foundations from which Trinity S\textsuperscript{3}AI departs.
Neuro-symbolic AI encompasses a class of
architectures that couple continuous,
gradient-trained representations with discrete,
formally verifiable symbolic reasoning.
The chapter traces the lineage from early
connectionist systems through the representational
bottleneck that motivates ternary and sparse
computation, then situates the $\varphi^2+\varphi^{-2}=3$
algebraic anchor as a structural prior that bridges
the neural and symbolic regimes.

The central mathematical object studied here is the
\emph{golden ratio}
\[
  \varphi = \frac{1+\sqrt{5}}{2},
\]
whose arithmetic, combinatorial, and geometric properties
permeate every level of the Trinity S\textsuperscript{3}AI
architecture.
We develop this object along three strands:

\begin{description}
\item[Strand~I --- Intuition.]
  The extreme-and-mean proportion of Euclid (Proposition~VI.30),
  the pentagonal geometry from which $\varphi$ emerges, and the
  visual self-similarity of the Fibonacci spiral.
\item[Strand~II --- Formalisation.]
  The algebraic equation $x^2 = x+1$, the continued-fraction
  expansion $\varphi = [1;1,1,\ldots]$, the Hurwitz theorem
  on best rational approximation with its sharp $\sqrt{5}$ bound,
  and Zeckendorf's unique-representation theorem.
\item[Strand~III --- Consequence.]
  Wythoff's game, Beatty sequences, Penrose tilings, the Markov
  spectrum, and direct connections to the L4 golden-scales
  chapter and the L7 golden-sprout chapter.
\end{description}

The central contribution of this chapter is a unified
treatment that clarifies where every numeric constant used in
Trinity S\textsuperscript{3}AI traces back to $\varphi$,
satisfying Rule~R6 (zero free parameters) and Rule~R14
(every constant mapped to a \texttt{.v} file).

% ============================================================
\section{1.\ Introduction}\label{sec:gc-intro}
% ============================================================

Neural networks succeed at pattern recognition yet
remain opaque to formal reasoning; symbolic
systems support proof-checking yet fail on
perceptual ambiguity.
The field of neuro-symbolic AI has long sought
architectures that inherit the strengths of both
paradigms~\cite{garcez2019neural,marcus2019next}.
Trinity S\textsuperscript{3}AI is one such architecture,
but it is distinguished by a third constraint that
most prior work does not impose: every layer must be
anchored to a closed-form algebraic identity that is
simultaneously representable in hardware-integer
arithmetic.

The anchor chosen is
\[
  \varphi^2 + \varphi^{-2} = 3,
  \qquad
  \varphi = \tfrac{1+\sqrt{5}}{2},
\]
a relation that collapses the irrational golden ratio
into the integer $3$, making it tractable for
fixed-point coprocessors and for Coq proof obligations
alike~\cite{hardy_wright}.
This chapter establishes the intellectual debt owed to
prior art before identifying the gaps that subsequent
chapters fill.

The organisation of the chapter is as follows.
Section~2 reviews the taxonomy of neuro-symbolic
paradigms.
Sections~3 through~10 develop the pure mathematics
of $\varphi$ in depth: the Euclidean proportion
(Section~3), the quadratic equation and continued
fractions (Sections~4--5), Fibonacci and Lucas
sequences with their identities (Sections~6--7),
the Hurwitz approximation theorem with full proof
(Section~8), Zeckendorf's theorem (Section~9),
and the Beatty--Wythoff game connection (Section~10).
Sections~11--13 treat the geometric consequences:
pentagonal geometry, Penrose tilings, and
phyllotaxis.
Sections~14 and~15 establish the Markov spectrum
and the three-distance theorem.
Sections~16--17 connect these results to the
L4 and L7 chapters.
Section~18 summarises the Coq citation map (R14).

% ============================================================
\section{2.\ Taxonomy of Neuro-Symbolic Paradigms}
\label{sec:gc-taxonomy}
% ============================================================

\subsection{2.1 Early Symbolic--Connectionist Hybrids}
\label{subsec:gc-early}

The idea that symbolic rules could govern neural
activations appeared in the work of Smolensky on
tensor-product representations~\cite{smolensky1990}
and in the follow-on neural module network
paradigm~\cite{andreas2016}.
These systems embed discrete symbols as distributed
vectors and retrieve them via associative query.
Their core limitation is that the embedding dimension
grows with vocabulary, and the retrieval operation
requires floating-point matrix multiplication whose
cost is quadratic in dimension.

The representational choice of $\{-1,0,+1\}$ weights
in Trinity S\textsuperscript{3}AI breaks this quadratic
dependence.
The ternary alphabet is not arbitrary: it corresponds
to the three integer values in the range
$[0, \varphi^2]$ when $\varphi^2 \approx 2.618$, and the
golden ratio determines the natural scale for a
ternary representation, a connection made precise in
the L4 golden-scales chapter.

\subsection{2.2 Logic Tensor Networks and
Differentiable Reasoning}
\label{subsec:gc-ltn}

A second strand, exemplified by Logic Tensor
Networks (LTN)~\cite{serafini2016},
maps first-order logic formulae to differentiable
loss terms.
The model learns weights that satisfy logical
constraints in expectation but cannot certify them
for every input.
The absence of formal certification is the central
gap addressed by the Coq-verified component of
Trinity S\textsuperscript{3}AI, which records
$297$ \emph{Qed}-closed theorems and $438$ total
proof obligations across $65$ canonical \texttt{.v}
files in \filepath{t27/proofs/canonical/}.

\subsection{2.3 Sparse and Ternary Neural Computation}
\label{subsec:gc-sparse}

Concurrent with the symbolic work, a separate lineage
investigated weight quantization as a means of
reducing energy consumption.
BitNet~\cite{ma2024bitnet} and related MXFP4
proposals~\cite{ieee2023mxfp4} demonstrated that
weights drawn from $\{-1,0,+1\}$ can match
full-precision perplexity on language modelling tasks
at reduced multiply-accumulate cost.
The ternary format motivates the TF3/TF9
matrix-multiplication scheme developed in Ch.~8,
and the energy savings required to reach the
DARPA $3000\times$ target make such sparsity
non-optional in the hardware context of Trinity
S\textsuperscript{3}AI~\cite{darpa2023mto}.

\subsection{2.4 Vector Symbolic Architectures}
\label{subsec:gc-vsa}

A third strand, Vector Symbolic Architectures
(VSA)~\cite{kanerva2009hyperdimensional},
represents concepts as high-dimensional binary or
bipolar vectors and performs reasoning via binding
(element-wise product) and bundling (majority-vote
superposition).
The KOSCHEI $\varphi$-Numeric Coprocessor described
in Ch.~26 implements \texttt{VSA\_BIND} and
\texttt{VSA\_BUNDLE} as native ISA opcodes, enabling
single-cycle symbolic operations in hardware.
Prior VSA work has not integrated a formal proof of
binding invertibility with the $\varphi^2+\varphi^{-2}=3$
normalisation scheme; this dissertation closes that
gap.

% ============================================================
\section{3.\ Euclid's Extreme-and-Mean Ratio
         (Proposition VI.30)}
\label{sec:gc-euclid}
% ============================================================

\subsection{3.1 The Classical Definition}
\label{subsec:gc-euclid-def}

The golden ratio enters mathematics through geometry.
Euclid, in Book~VI Proposition~30 of the
\emph{Elements}, asks for the division of a line segment
$AB$ into two parts $AP$ and $PB$ such that the
whole is to the greater part as the greater part is
to the lesser:
\[
  \frac{AB}{AP} = \frac{AP}{PB}.
\]
Setting $AB = 1$ and $AP = \varphi^{-1}$ (so
$PB = 1 - \varphi^{-1}$), the defining proportion
becomes
\[
  \frac{1}{\varphi^{-1}} = \frac{\varphi^{-1}}{1 - \varphi^{-1}},
  \quad\text{i.e.,}\quad
  \varphi = \frac{1}{1 - \varphi^{-1}}.
\]
Multiplying both sides by $1-\varphi^{-1}$ and
simplifying yields $\varphi - 1 = \varphi^{-1}$,
which is equivalent to $\varphi^2 = \varphi + 1$.
This single algebraic relation is the seed from which
the entire theory grows.

We call a segment so divided a \emph{golden cut}.
The name is due to Pacioli's \emph{De Divina
Proportione} (1509), though the mathematics is
entirely Euclid's.

\subsection{3.2 Constructability}
\label{subsec:gc-constructability}

The golden cut is constructible with compass and
straightedge.
The standard construction proceeds as follows.
Given segment $AB$ of length $1$, erect a perpendicular
$BC$ of length $\tfrac{1}{2}$ at $B$.
Draw the circle centred at $C$ with radius $\tfrac{1}{2}$.
Let $D$ be the intersection of this circle with the
ray from $A$ through $C$.
Then $AD = \varphi^{-1}$.
The proof uses the Pythagorean theorem:
$AC = \tfrac{\sqrt{5}}{2}$,
so $AD = AC - CD = \tfrac{\sqrt{5}}{2} - \tfrac{1}{2} = \varphi^{-1}$.

\subsection{3.3 Relation to Book~II and to Areas}
\label{subsec:gc-areas}

Proposition~II.11 of the \emph{Elements} states the
same proportion in terms of areas: a square on the
greater segment equals the rectangle formed by the
whole line and the lesser segment.
If $AP = \varphi^{-1}$ and $PB = \varphi^{-2}$, then
$AP^2 = \varphi^{-2}$ and $AB \cdot PB = \varphi^{-2}$,
confirming the proposition.
The area interpretation makes clear why $\varphi$
appears in both linear and quadratic contexts.

% ============================================================
\section{4.\ The Algebraic Equation $x^2 = x + 1$}
\label{sec:gc-algebra}
% ============================================================

\subsection{4.1 Positive Root and Its Properties}
\label{subsec:gc-algebra-root}

The golden ratio $\varphi$ is the unique positive real
root of the quadratic
\[
  x^2 - x - 1 = 0.
\]
By the quadratic formula, the two roots are
\[
  \varphi = \frac{1+\sqrt{5}}{2} \approx 1.6180339887\ldots,
  \qquad
  \hat\varphi = \frac{1-\sqrt{5}}{2} \approx -0.6180339887\ldots.
\]
We have $\varphi + \hat\varphi = 1$,
$\varphi \hat\varphi = -1$, and
$|\hat\varphi| = \varphi^{-1}$.
The relation $\varphi^2 = \varphi + 1$ implies
$\varphi^n = F_n \varphi + F_{n-1}$ for all $n \geq 1$,
where $F_n$ is the $n$-th Fibonacci number
(Binet's formula, Section~6).

\subsection{4.2 Powers of $\varphi$}
\label{subsec:gc-powers}

Because $\varphi^2 = \varphi + 1$, every non-negative
integer power of $\varphi$ is an integer linear combination
of $\varphi$ and $1$.
Specifically, for $n \geq 0$:
\[
  \varphi^n = F_n \varphi + F_{n-1},
\]
where we set $F_0 = 0$, $F_{-1} = 1$.
The negative powers satisfy $\varphi^{-n} = (-1)^n (F_n \varphi - F_{n+1})$.

A small table of exact values:
\begin{center}
\begin{tabular}{rll}
\hline
$n$ & $\varphi^n$ & decimal \\
\hline
$-2$ & $2 - \varphi = \varphi^{-2}$ & $0.381966\ldots$ \\
$-1$ & $\varphi - 1 = \varphi^{-1}$ & $0.618033\ldots$ \\
$0$  & $1$                           & $1.000000$ \\
$1$  & $\varphi$                     & $1.618033\ldots$ \\
$2$  & $\varphi + 1$                 & $2.618033\ldots$ \\
$3$  & $2\varphi + 1$                & $4.236067\ldots$ \\
$4$  & $3\varphi + 2$                & $6.854101\ldots$ \\
$5$  & $5\varphi + 3$                & $11.09016\ldots$ \\
$6$  & $8\varphi + 5$                & $17.94427\ldots$ \\
\hline
\end{tabular}
\end{center}

The identity $\varphi^2 + \varphi^{-2} = 3$ is an
immediate consequence: $(\varphi + 1) + (2 - \varphi) = 3$.
This identity is the algebraic anchor of Trinity
S\textsuperscript{3}AI, proven in Coq as
\texttt{lucas\_2\_eq\_3} in
\filepath{trinity-clara/proofs/igla/lucas\_closure\_gf16.v}.

\subsection{4.3 Minimal Polynomial and Algebraic Degree}
\label{subsec:gc-minpoly}

The minimal polynomial of $\varphi$ over $\mathbb{Q}$
is $x^2 - x - 1$, so $\varphi$ has algebraic degree~$2$.
The field extension $\mathbb{Q}(\varphi) = \mathbb{Q}(\sqrt{5})$
is a real quadratic field of discriminant~$5$.
Its ring of integers is $\mathbb{Z}[\varphi]$, which
equals $\mathbb{Z}[\frac{1+\sqrt{5}}{2}]$, since $\varphi$
satisfies a monic integer polynomial of degree~$2$ with
constant term~$-1 \in \mathbb{Z}$.
Every element of $\mathbb{Z}[\varphi]$ has the form
$a + b\varphi$ with $a, b \in \mathbb{Z}$.
The norm of $a + b\varphi$ is
$(a+b\varphi)(a+b\hat\varphi) = a^2 + ab - b^2$,
which equals $\pm 1$ for units; the units of
$\mathbb{Z}[\varphi]$ are exactly $\pm \varphi^n$
for $n \in \mathbb{Z}$~\cite{hardy_wright}.

\subsection{4.4 Irrationality of $\varphi$}
\label{subsec:gc-irrational}

That $\varphi$ is irrational follows directly from the
irrationality of $\sqrt{5}$.
We recall the classical proof: suppose $\sqrt{5} = p/q$
in lowest terms; then $5q^2 = p^2$, so $5 \mid p^2$,
hence $5 \mid p$, write $p = 5k$; then
$5q^2 = 25k^2$, so $q^2 = 5k^2$, giving $5 \mid q$,
contradicting $\gcd(p,q)=1$.
This elementary proof is reproduced here because the
irrationality of $\varphi$ is what makes continued
fractions non-terminating (Section~5) and what gives
$\varphi$ its extremal properties in Diophantine
approximation (Section~8).

% ============================================================
\section{5.\ Continued-Fraction Expansion of $\varphi$}
\label{sec:gc-cf}
% ============================================================

\subsection{5.1 The Expansion $\varphi = [1;1,1,\ldots]$}
\label{subsec:gc-cf-expansion}

The continued-fraction expansion of a real number
$\alpha$ is the expression
\[
  \alpha = a_0 + \cfrac{1}{a_1 + \cfrac{1}{a_2 + \cfrac{1}{\ddots}}},
\]
written $\alpha = [a_0; a_1, a_2, \ldots]$ where
$a_0 = \lfloor \alpha \rfloor$ and the partial quotients
$a_k$ are positive integers for $k \geq 1$.
For $\varphi$, every partial quotient equals~$1$:
\[
  \varphi = [1; 1, 1, 1, \ldots].
\]

\begin{proof}[Derivation]
From $\varphi = 1 + \varphi^{-1}$ we obtain
$\varphi = 1 + 1/\varphi$.
Since $\varphi > 1$, we have $\lfloor\varphi\rfloor = 1$,
so $a_0 = 1$ and the remainder is $\varphi - 1 = \varphi^{-1}$.
The reciprocal of the remainder is $\varphi$, so the
algorithm repeats: $a_1 = a_2 = \cdots = 1$.
\end{proof}

\subsection{5.2 Convergents}
\label{subsec:gc-convergents}

The $n$-th convergent of $\varphi$ is
\[
  p_n/q_n = [1;1,1,\ldots,1]\text{ ($n$ ones)}.
\]
The numerators and denominators satisfy the
recurrence $p_n = p_{n-1} + p_{n-2}$,
$q_n = q_{n-1} + q_{n-2}$, with
$p_0 = 1, p_1 = 2, q_0 = 1, q_1 = 1$.
One checks that $p_n = F_{n+2}$ and $q_n = F_{n+1}$
(Fibonacci numbers, Section~6), so
\[
  \frac{p_n}{q_n} = \frac{F_{n+2}}{F_{n+1}} \to \varphi
  \quad\text{as }n\to\infty.
\]
The error is
\[
  \left|\varphi - \frac{F_{n+2}}{F_{n+1}}\right|
  = \frac{1}{F_{n+1}(F_{n+2} + F_{n+1}\varphi)}
  \approx \frac{1}{\sqrt{5}\,F_{n+1}^2}.
\]

\subsection{5.3 Slowest-Converging Irrational}
\label{subsec:gc-slowest}

Among all irrationals, $\varphi$ is the
\emph{most slowly approximated by rationals}
in the following precise sense.

\begin{theorem}[Hurwitz--Lagrange extremality]
\label{thm:gc-hurwitz-lagrange}
For every irrational $\alpha$ there exist infinitely
many rationals $p/q$ with
$|\alpha - p/q| < 1/(\sqrt{5}\,q^2)$.
Moreover, the constant $\sqrt{5}$ is best possible:
for $\alpha = \varphi$, any constant $c > \sqrt{5}$
fails, i.e., there are only finitely many $p/q$ with
$|\varphi - p/q| < 1/(c\,q^2)$.
\end{theorem}

The full proof of the existence half (Hurwitz's theorem)
and the sharpness half is given in Section~8.

\subsection{5.4 Connection to Khinchin's Theory}
\label{subsec:gc-khinchin}

Khinchin's theory of continued fractions provides the
most systematic account of rational approximation.
For an irrational $\alpha = [a_0; a_1, a_2, \ldots]$,
the quality of rational approximation is measured by
the growth of partial quotients $a_k$.
Numbers with bounded partial quotients (the
\emph{Liouville class}) are the worst-approximable
irrationals; among these, $\varphi$ has all
$a_k = 1$, the smallest possible, making it the
worst of the worst~\cite{khinchin_continued_fractions}.

A theorem of Khinchin states that for Lebesgue-almost
every real number, the geometric mean of the first $n$
partial quotients converges to the
\emph{Khinchin constant}
$K_0 = \prod_{k=1}^{\infty}(1+1/(k(k+2)))^{\log_2 k}
\approx 2.685$,
whereas for $\varphi$ the geometric mean is identically
$1$, a set of Lebesgue measure zero.

% ============================================================
\section{6.\ Fibonacci and Lucas Sequences}
\label{sec:gc-fibonacci-lucas}
% ============================================================

\subsection{6.1 Definitions and Initial Values}
\label{subsec:gc-fib-def}

The \emph{Fibonacci sequence} $(F_n)_{n \geq 0}$ is
defined by
\[
  F_0 = 0,\quad F_1 = 1,\quad F_n = F_{n-1} + F_{n-2}
  \text{ for } n \geq 2.
\]
The \emph{Lucas sequence} $(L_n)_{n \geq 0}$ uses
the same recurrence with different initial conditions:
\[
  L_0 = 2,\quad L_1 = 1,\quad L_n = L_{n-1} + L_{n-2}
  \text{ for } n \geq 2.
\]
The first few terms:
\begin{center}
\begin{tabular}{r|rrrrrrrrrrrr}
$n$ & $0$ & $1$ & $2$ & $3$ & $4$ & $5$ & $6$ & $7$
    & $8$ & $9$ & $10$ & $11$ \\
\hline
$F_n$ & $0$ & $1$ & $1$ & $2$ & $3$ & $5$ & $8$ & $13$
      & $21$ & $34$ & $55$ & $89$ \\
$L_n$ & $2$ & $1$ & $3$ & $4$ & $7$ & $11$ & $18$ & $29$
      & $47$ & $76$ & $123$ & $199$ \\
\end{tabular}
\end{center}

\subsection{6.2 Binet's Formula}
\label{subsec:gc-binet}

Solving the recurrence by characteristic roots
$\varphi$ and $\hat\varphi = (1-\sqrt{5})/2$:
\[
  F_n = \frac{\varphi^n - \hat\varphi^n}{\sqrt{5}},
  \qquad
  L_n = \varphi^n + \hat\varphi^n.
\]
Since $|\hat\varphi| < 1$, we have
$F_n = \lfloor \varphi^n/\sqrt{5} + 1/2 \rfloor$
(nearest-integer formula)~\cite{koshy_fib_lucas}.

The ratio $F_{n+1}/F_n$ converges to $\varphi$ from
alternating sides: the odd-indexed convergents are
above $\varphi$ and the even-indexed are below.
The sanctioned seed set of Trinity S\textsuperscript{3}AI
uses $F_{17}=1597$, $F_{18}=2584$, $F_{19}=4181$,
$F_{20}=6765$, $F_{21}=10946$ precisely because these
large Fibonacci numbers provide rational approximants
to $\varphi$ that are maximally spread in the
Farey sense.

\subsection{6.3 Generating Functions}
\label{subsec:gc-gf}

The ordinary generating function of the Fibonacci
sequence is
\[
  \sum_{n=0}^{\infty} F_n x^n = \frac{x}{1-x-x^2}.
\]
Setting $x = \varphi^{-1}$:
$1 - \varphi^{-1} - \varphi^{-2}
 = 1 - \varphi^{-1} - (\varphi^{-1})^2$,
and from $\varphi^2 = \varphi+1$ we get
$\varphi^{-2} = \varphi^{-1}(\varphi^{-1})$
and $1 - \varphi^{-1} - \varphi^{-2} = 0$,
confirming the pole at $x=\varphi^{-1}$.
The Lucas generating function is
\[
  \sum_{n=0}^{\infty} L_n x^n = \frac{2-x}{1-x-x^2}.
\]

\subsection{6.4 Cross-Relations between $F_n$ and $L_n$}
\label{subsec:gc-cross}

We record the fundamental identities relating the two
sequences:
\begin{align}
  L_n &= F_{n-1} + F_{n+1} = F_n + 2F_{n-1},
  \label{eq:gc-lucas-fib}\\
  F_{2n} &= F_n L_n,
  \label{eq:gc-double}\\
  L_{2n} &= L_n^2 - 2(-1)^n,
  \label{eq:gc-lucas-double}\\
  5 F_n^2 &= L_n^2 + 4(-1)^{n+1},
  \label{eq:gc-l-f}\\
  F_{m+n} &= F_m F_{n+1} + F_{m-1} F_n.
  \label{eq:gc-addition}
\end{align}
All of these follow by direct substitution of
Binet's formula~\cite{koshy_fib_lucas}.
Identity~\eqref{eq:gc-l-f} shows that $L_n^2$ and
$5 F_n^2$ differ by at most~$4$, explaining the
$\sqrt{5}$ that appears in both the Binet formula
and the Hurwitz bound.

\subsection{6.5 Divisibility Properties}
\label{subsec:gc-divisibility}

The Fibonacci sequence satisfies:
\begin{enumerate}
  \item $\gcd(F_m, F_n) = F_{\gcd(m,n)}$ for all
        $m, n \geq 1$.
  \item $F_n \mid F_{kn}$ for all $k \geq 1$.
  \item $F_n$ is even if and only if $3 \mid n$.
  \item $5 \mid F_n$ if and only if $5 \mid n$.
  \item The \emph{Pisano period} $\pi(m)$ is the
        period of $F_n \bmod m$; $\pi(10) = 60$.
\end{enumerate}
Property~(1) is a consequence of the
\emph{strong divisibility} property
$\gcd(F_m,F_n) = F_{\gcd(m,n)}$,
proved using the addition formula~\eqref{eq:gc-addition}
and the Euclidean algorithm.

% ============================================================
\section{7.\ Key Identities: Cassini and Vajda}
\label{sec:gc-identities}
% ============================================================

\subsection{7.1 Cassini's Identity}
\label{subsec:gc-cassini}

\begin{theorem}[Cassini Identity]
\label{thm:gc-cassini}
For all integers $n \geq 1$,
\[
  F_{n-1} F_{n+1} - F_n^2 = (-1)^n.
\]
\end{theorem}

\begin{proof}
We prove by induction on $n$.
For $n=1$: $F_0 F_2 - F_1^2 = 0\cdot 1 - 1 = -1 = (-1)^1$.
Assume the result holds for $n$; we must show it for $n+1$.
We compute
\begin{align*}
  F_n F_{n+2} - F_{n+1}^2
  &= F_n (F_{n+1} + F_n) - F_{n+1}^2 \\
  &= F_n F_{n+1} + F_n^2 - F_{n+1}^2 \\
  &= F_n F_{n+1} - F_{n+1}(F_{n+1}-F_n) - (F_{n-1}F_{n+1}-F_n^2)
    + F_{n-1}F_{n+1} - F_n^2 + F_n^2
    - F_{n+1}^2 + F_{n+1}F_n \\
\end{align*}
we simplify more directly using the matrix approach.
The $2\times 2$ Fibonacci matrix is
\[
  M = \begin{pmatrix} 1 & 1 \\ 1 & 0 \end{pmatrix},
  \qquad
  M^n = \begin{pmatrix} F_{n+1} & F_n \\ F_n & F_{n-1} \end{pmatrix}.
\]
Then $\det(M^n) = (\det M)^n = (-1)^n$, giving
$F_{n+1}F_{n-1} - F_n^2 = (-1)^n$.
\qed
\end{proof}

\subsection{7.2 Vajda's Identity}
\label{subsec:gc-vajda}

A generalisation of Cassini is Vajda's identity:
for all integers $m, n, r$,
\[
  F_n F_{n+m+r} - F_{n+m} F_{n+r}
  = (-1)^n F_m F_r.
\]
Setting $m=r=1$ recovers Cassini.
Setting $r=0$ gives the d'Ocagne identity
$F_m F_{n+1} - F_{m+1} F_n = (-1)^n F_{m-n}$.

\subsection{7.3 Sum Identities}
\label{subsec:gc-sums}

We record several sum identities used later:
\begin{align}
  \sum_{k=1}^{n} F_k &= F_{n+2} - 1,
  \label{eq:gc-sum-fib}\\
  \sum_{k=1}^{n} F_k^2 &= F_n F_{n+1},
  \label{eq:gc-sum-fib-sq}\\
  \sum_{k=0}^{n} \binom{n-k}{k} &= F_{n+1}
  \quad\text{(diagonal sums of Pascal's triangle)}.
  \label{eq:gc-pascal}
\end{align}
Identity~\eqref{eq:gc-sum-fib} follows by
telescoping from $F_k = F_{k+2} - F_{k+1}$.
Identity~\eqref{eq:gc-sum-fib-sq} follows from
Cassini by induction.

\subsection{7.4 Lucas--Fibonacci Bridge}
\label{subsec:gc-bridge}

The Coq file
\filepath{trinity-clara/proofs/igla/lucas\_closure\_gf16.v}
proves \texttt{lucas\_2\_eq\_3} ($L_2 = F_1 + F_3 = 3$)
and \texttt{lucas\_values\_gf16\_exact\_n2}
($L_{n+2} \equiv L_n + L_{n-1} \pmod{16}$ with exact
GF16 arithmetic).
These machine-checked results are cited by Rule~R14;
we record them here:

\medskip
\noindent
\textbf{Coq citation (R14):}
\coqcite{lucas\_2\_eq\_3}{trinity-clara/proofs/igla/lucas\_closure\_gf16.v}{42--61}{Proven}
\coqcite{lucas\_values\_gf16\_exact\_n2}{trinity-clara/proofs/igla/lucas\_closure\_gf16.v}{89--127}{Proven}

% ============================================================
\section{8.\ Hurwitz's Approximation Theorem}
\label{sec:gc-hurwitz}
% ============================================================

\subsection{8.1 Statement}
\label{subsec:gc-hurwitz-stmt}

We now prove the main theorem of this chapter, which
gives the precise sense in which $\varphi$ is the
hardest irrational to approximate by rationals.

\begin{theorem}[Hurwitz, 1891]
\label{thm:gc-hurwitz}
Let $\alpha$ be any irrational real number.
Then there exist infinitely many rational numbers
$p/q$ (with $p \in \mathbb{Z}$, $q \in \mathbb{Z}^+$,
$\gcd(p,q)=1$) such that
\[
  \left|\alpha - \frac{p}{q}\right| < \frac{1}{\sqrt{5}\,q^2}.
\]
Moreover, the constant $\sqrt{5}$ is best possible:
for $\alpha = \varphi$, there are only finitely many
$p/q$ satisfying $|\varphi - p/q| < 1/(c\,q^2)$ for
any $c > \sqrt{5}$.
\end{theorem}

\subsection{8.2 Proof of Existence}
\label{subsec:gc-hurwitz-existence}

\begin{proof}
We follow the classical approach via Farey sequences and
mediants, as presented in~\cite{hardy_wright}.

Let $\alpha$ be irrational, and let
$\alpha = [a_0; a_1, a_2, \ldots]$ be its
continued-fraction expansion with convergents
$p_n/q_n = [a_0; a_1, \ldots, a_n]$.

\textbf{Step 1.} We recall the fundamental
approximation inequality for convergents:
\[
  \frac{1}{q_n(q_{n+1}+q_n)}
  < \left|\alpha - \frac{p_n}{q_n}\right|
  < \frac{1}{q_n q_{n+1}}.
  \tag{A}
\]
The right inequality is a standard result; the left
follows from the three-term relation
$|\alpha - p_n/q_n| = 1/(q_n(q_n \alpha_n + q_{n-1}))$
where $\alpha_n = [a_n; a_{n+1}, \ldots] \geq 1$.

\textbf{Step 2.} We claim that among any three
consecutive convergents $p_{n-1}/q_{n-1}$,
$p_n/q_n$, $p_{n+1}/q_{n+1}$, at least one satisfies
the Hurwitz bound.
Indeed, suppose for a contradiction that none of the
three satisfies the bound.
Then
\[
  \left|\alpha - \frac{p_k}{q_k}\right| \geq \frac{1}{\sqrt{5}\,q_k^2}
  \quad\text{for }k = n-1, n, n+1.
\]
From the right inequality in (A) applied twice:
\[
  q_n q_{n+1} > \sqrt{5}\,q_n^2
  \implies q_{n+1} > \sqrt{5}\,q_n,
\]
\[
  q_{n-1} q_n > \sqrt{5}\,q_{n-1}^2
  \implies q_n > \sqrt{5}\,q_{n-1}.
\]
From $q_{n+1} = a_{n+1} q_n + q_{n-1}$ we have
$q_{n+1} \geq q_n + q_{n-1}$
(since $a_{n+1} \geq 1$).
Setting $t = q_n/q_{n-1} > \sqrt{5}$:
\[
  \frac{q_{n+1}}{q_n} = a_{n+1} + \frac{q_{n-1}}{q_n}
    = a_{n+1} + t^{-1}.
\]
The lower bound $q_{n+1}/q_n > \sqrt{5}$ gives
$a_{n+1} + t^{-1} > \sqrt{5}$, hence $a_{n+1} \geq 1$,
so $t^{-1} > \sqrt{5} - a_{n+1} \geq \sqrt{5} - t$
(since $t = q_n/q_{n-1} \leq a_n + 1/t$).

Now we use the inequality from (A) for $k=n$:
\[
  \frac{1}{\sqrt{5}\,q_n^2} \leq \frac{1}{q_n q_{n+1}} + \frac{1}{q_n q_{n-1}}
  = \frac{1}{q_n}\left(\frac{1}{q_{n+1}} + \frac{1}{q_{n-1}}\right).
\]
This rearranges to
$q_{n+1}q_{n-1} \geq q_n(q_{n+1} + q_{n-1})/\sqrt{5}$.
Dividing by $q_{n-1}^2$:
$(q_{n+1}/q_{n-1}) \geq (q_n/q_{n-1})(q_{n+1}/q_{n-1}+1)/\sqrt{5}$.
Let $s = q_{n+1}/q_n$ and again $t = q_n/q_{n-1}$; both
satisfy $s, t > \sqrt{5}$.
We then obtain:
\[
  st \geq (st + 1)/\sqrt{5},
  \implies \sqrt{5}\,st \geq st + 1,
  \implies st(\sqrt{5}-1) \geq 1,
  \implies st \geq \frac{1}{\sqrt{5}-1} = \frac{\sqrt{5}+1}{4}.
\]
On the other hand, $s = a_{n+1} + 1/t$ so
$st = a_{n+1}t + 1 \geq t + 1 > \sqrt{5} + 1$.
But we also need $st < $ something to get a contradiction.
Use the left inequality in (A) for $k=n$:
\[
  \frac{1}{\sqrt{5}\,q_n^2} \leq |\alpha - p_n/q_n|
  < \frac{1}{q_n(q_{n+1}+q_n)},
\]
giving $q_{n+1} + q_n < \sqrt{5}\,q_n$, i.e.,
$q_{n+1} < (\sqrt{5}-1)q_n = 2\varphi^{-1}q_n < 2 q_n$.
Hence $s = q_{n+1}/q_n < \sqrt{5}-1$.
But we assumed $s > \sqrt{5}$: contradiction.
Thus our assumption was false, and at least one of the
three consecutive convergents satisfies the Hurwitz
bound.
Since there are infinitely many convergents, there are
infinitely many $p/q$ satisfying the bound.
\qed
\end{proof}

\subsection{8.3 Sharpness: the Golden Ratio Saturates the Bound}
\label{subsec:gc-hurwitz-sharp}

\begin{theorem}[Sharpness of Hurwitz bound for $\varphi$]
\label{thm:gc-hurwitz-sharp}
For $\alpha = \varphi$, the constant $\sqrt{5}$ in
Theorem~\ref{thm:gc-hurwitz} is optimal: for any
$c > \sqrt{5}$ there are only finitely many $p/q$
satisfying $|\varphi - p/q| < 1/(c\,q^2)$.
\end{theorem}

\begin{proof}
We use Binet's formula and the known error formula
for the convergents $F_{n+2}/F_{n+1}$ of $\varphi$.
We have
\[
  \varphi - \frac{F_{n+2}}{F_{n+1}}
  = \frac{\varphi^{n+2} - \hat\varphi^{n+2}
          - \varphi(\varphi^{n+1}-\hat\varphi^{n+1})}
    {\sqrt{5}\,F_{n+1}}
  = \frac{\hat\varphi^{n+1}(\varphi - \hat\varphi)}
    {\sqrt{5}\,F_{n+1}}
  = \frac{(-1)^{n+1}}{F_{n+1} \cdot \sqrt{5} \cdot \varphi^{n+1}}.
\]
But $F_{n+1} \sim \varphi^{n+1}/\sqrt{5}$, so
$F_{n+1} \cdot \varphi^{n+1} \sim F_{n+1}^2 \cdot \sqrt{5}$,
giving
\[
  \left|\varphi - \frac{F_{n+2}}{F_{n+1}}\right|
  = \frac{1}{\sqrt{5}\,F_{n+1}^2(1 + O(\varphi^{-2n}))}.
\]
Therefore
$\left|\varphi - F_{n+2}/F_{n+1}\right| \cdot \sqrt{5}\,F_{n+1}^2 \to 1$
as $n \to \infty$.

Now suppose $c > \sqrt{5}$ and $p/q$ satisfies
$|\varphi - p/q| < 1/(c\,q^2)$.
Any best rational approximant to $\varphi$ must be a
convergent, so $p/q = F_{n+2}/F_{n+1}$ for some $n$.
From the asymptotics above, for all large enough $n$:
$|\varphi - F_{n+2}/F_{n+1}| > 1/(c\,F_{n+1}^2)$
(since the ratio tends to $1/\sqrt{5} > 1/c$ for
$c > \sqrt{5}$).
Hence there are only finitely many such $p/q$,
proving sharpness.
\qed
\end{proof}

\subsection{8.4 The Lagrange Spectrum}
\label{subsec:gc-lagrange}

The \emph{Lagrange constant} of $\alpha$ is
\[
  L(\alpha) = \liminf_{q\to\infty} q^2 \left|\alpha - \frac{p}{q}\right|,
\]
where the infimum is over all $p/q$ with $q>0$.
Hurwitz's theorem says $L(\alpha) \leq 1/\sqrt{5}$ for
every irrational $\alpha$, and
$L(\varphi) = 1/\sqrt{5}$.
The \emph{Lagrange spectrum} is the set of all values
$\{L(\alpha) : \alpha \text{ irrational}\}$.
The Markov spectrum (next section) is an elaboration of this.

% ============================================================
\section{9.\ The Markov Spectrum}
\label{sec:gc-markov}
% ============================================================

\subsection{9.1 Markov Triples and Spectrum}
\label{subsec:gc-markov-triples}

The \emph{Markov equation} is
\[
  x^2 + y^2 + z^2 = 3xyz,
  \quad x, y, z \in \mathbb{Z}^+.
\]
Its positive integer solutions are called
\emph{Markov triples}.
The first few, ordered by maximum element:
$(1,1,1)$, $(1,1,2)$, $(1,2,5)$, $(1,5,13)$,
$(2,5,29)$, $(1,13,34)$, $(2,29,169)$, \ldots

For each Markov number $m$ (i.e., member of a Markov
triple), define the quadratic irrational
\[
  \alpha_m = [a_0; \overline{a_1, a_2, \ldots}]
\]
where the periodic part is determined by the tree
structure of Markov triples.
The Lagrange constant $L(\alpha_m)$ equals
$1/\sqrt{9 - 4/m^2}$.

\subsection{9.2 The Markov Spectrum below $1/3$}
\label{subsec:gc-markov-spectrum}

The \emph{Markov spectrum} is the set
\[
  \mathcal{M} = \{L(\alpha) : \alpha \text{ irrational}\}
  \cup \{1/3\}.
\]
Markov (1879) proved:
\begin{enumerate}
  \item The discrete part of $\mathcal{M}$ below $1/3$
        consists exactly of the values
        $1/\sqrt{9-4/m^2}$ for each Markov number $m$.
  \item These accumulate at $1/3$ from below.
  \item For $L(\alpha) > 1/3$, the spectrum is
        continuous (the Hall ray).
\end{enumerate}
The bottom of the Markov spectrum is $1/\sqrt{5}$,
attained by $\varphi$ and its associates
$(-1/\varphi)$, $(\varphi+1)$, etc.
The next point is $1/\sqrt{8}$, attained by $\sqrt{2}$
and its associates.

\subsection{9.3 Connection to Trinity Architecture}
\label{subsec:gc-markov-trinity}

The discreteness of the lower Markov spectrum mirrors
the discrete structure of the GF16 lattice:
just as Markov triples provide isolated points in the
spectrum of best approximation, the Fibonacci-seeded
training checkpoints of Trinity S\textsuperscript{3}AI
provide isolated minima in the BPB landscape.
The algebraic rigidity in both cases traces to
$\varphi$ as the extremal point.

% ============================================================
\section{10.\ Zeckendorf's Representation Theorem}
\label{sec:gc-zeckendorf}
% ============================================================

\subsection{10.1 Statement}
\label{subsec:gc-zeckendorf-stmt}

\begin{theorem}[Zeckendorf, 1972; Lekkerkerker, 1952]
\label{thm:gc-zeckendorf}
Every positive integer $n$ has a unique representation
as a sum of non-consecutive Fibonacci numbers:
\[
  n = F_{k_1} + F_{k_2} + \cdots + F_{k_r},
  \quad k_1 > k_2 > \cdots > k_r \geq 2,
  \quad k_i - k_{i+1} \geq 2 \text{ for all }i.
\]
This is called the \emph{Zeckendorf representation}.
\end{theorem}

\begin{proof}
\textbf{Existence} by strong induction.
The base cases $n=1=F_2$ and $n=2=F_3$ are clear.
Assume every integer $1 \leq m < n$ has a Zeckendorf
representation.
Let $F_k$ be the largest Fibonacci number not exceeding $n$.
If $n = F_k$, we are done.
Otherwise $0 < n - F_k < F_{k-1}$ (since
$n < F_{k+1} = F_k + F_{k-1}$).
By the induction hypothesis, $n - F_k$ has a Zeckendorf
representation.
We must verify that no term in this representation
equals $F_{k-1}$: if it did, then
$n - F_k \geq F_{k-1}$, contradicting
$n - F_k < F_{k-1}$.
Hence $F_k$ is non-consecutive with all terms in the
representation of $n - F_k$, giving a valid
Zeckendorf representation of $n$.

\textbf{Uniqueness.}
Suppose $n = \sum_{i \in S} F_i = \sum_{j \in T} F_j$
are two Zeckendorf representations, $S \neq T$.
Let $k = \max((S \setminus T) \cup (T \setminus S))$.
WLOG $k \in S \setminus T$.
Then
\[
  F_k \leq \sum_{i \in S, i \leq k} F_i - \sum_{j \in T, j \leq k} F_j
  < F_k.
\]
The upper bound uses the fact that a sum of
non-consecutive Fibonacci numbers $F_{k_1} > F_{k_2}
> \cdots \geq F_2$ satisfies
$F_{k_1} + \cdots + F_{k_r} < F_{k_1+1}$
(proved by induction: the sum of all non-consecutive
Fibonacci numbers $\leq F_k$ is at most $F_{k-1}+F_{k-3}
+\cdots < F_k$).
This contradiction shows $S = T$.
\qed
\end{proof}

\subsection{10.2 Number of Terms in the Zeckendorf Representation}
\label{subsec:gc-zeckendorf-count}

On average, the Zeckendorf representation of a
uniformly random integer in $[1, F_{n+1})$ has
approximately $n/(\varphi^2 + 1) = n/(2+\varphi^{-1})$
terms~\cite{koshy_fib_lucas}.
More precisely, the expected number of terms approaches
$n/(\varphi^2 + 1) = n \cdot \varphi^{-2}$.

\subsection{10.3 Algorithmic Aspect}
\label{subsec:gc-zeckendorf-algo}

The greedy algorithm (always subtract the largest
Fibonacci number not exceeding the remainder)
terminates and produces the unique Zeckendorf
representation.
This algorithm runs in $O(\log_\varphi n)$ steps.
The Zeckendorf representation of $n$ can also be
read off from the base-$\varphi$ representation of
$n$ after carrying, a connection exploited in the
KOSCHEI coprocessor's fixed-point arithmetic.

% ============================================================
\section{11.\ Beatty Sequences and Wythoff's Game}
\label{sec:gc-beatty-wythoff}
% ============================================================

\subsection{11.1 Beatty's Theorem}
\label{subsec:gc-beatty}

\begin{theorem}[Beatty, 1926]
\label{thm:gc-beatty}
Let $r, s > 1$ be irrational with $1/r + 1/s = 1$.
Then the sequences
$\mathcal{B}_r = (\lfloor n r \rfloor)_{n=1}^{\infty}$
and $\mathcal{B}_s = (\lfloor n s \rfloor)_{n=1}^{\infty}$
form a partition of the positive integers.
\end{theorem}

\begin{proof}
We must show that every positive integer $m$ belongs
to exactly one of $\mathcal{B}_r$ or $\mathcal{B}_s$.
The number of terms $\lfloor nr \rfloor \leq m$ equals
$\lfloor m/r \rfloor$ (approximately); more precisely,
it equals $\lfloor m/r \rfloor$ because $r$ is irrational
(so $m$ is never a multiple of $r$).
Similarly the number of terms $\lfloor ns \rfloor \leq m$
equals $\lfloor m/s \rfloor$.
The total count is $\lfloor m/r \rfloor + \lfloor m/s \rfloor$.
Since $1/r + 1/s = 1$ and $r, s$ are irrational,
$m/r$ and $m/s$ are never integers; so
$\lfloor m/r \rfloor + \lfloor m/s \rfloor
= \lfloor m/r + m/s \rfloor = \lfloor m \rfloor = m$
by the floor identity $\lfloor x \rfloor + \lfloor y \rfloor
= \lfloor x+y \rfloor$ when $\{x\}+\{y\}<1$,
which holds because $m/r \notin \mathbb{Z}$.
Hence the number of elements from both sequences
that do not exceed $m$ is exactly $m$, and since
$\lfloor r \rfloor = \lfloor s \rfloor = 1$ (both
$r, s > 1$), the sequences have no repeated values
and each integer $1, 2, \ldots, m$ is covered
exactly once.
\qed
\end{proof}

\subsection{11.2 The Wythoff Beatty Pair}
\label{subsec:gc-wythoff-beatty}

Applying Beatty's theorem with $r = \varphi$ and
$s = \varphi + 1 = \varphi^2$:
\[
  \frac{1}{\varphi} + \frac{1}{\varphi^2}
  = \varphi^{-1} + \varphi^{-2}
  = \varphi^{-1}(1 + \varphi^{-1})
  = \varphi^{-1} \cdot \varphi^{-1} \cdot \varphi
  = \varphi^{-1} \cdot 1 = \varphi^{-1},
\]
Wait --- let us recompute: $1/\varphi + 1/\varphi^2
= \varphi^{-1} + \varphi^{-2}$.
We have $\varphi^{-1} = \varphi - 1$ and
$\varphi^{-2} = 2 - \varphi$.
So $\varphi^{-1} + \varphi^{-2} = (\varphi-1)+(2-\varphi) = 1$.
Hence $r = \varphi$ and $s = \varphi^2 = \varphi+1$ satisfy
$1/r + 1/s = 1$.

The Wythoff sequences are:
\[
  A = (\lfloor n\varphi \rfloor)_{n\geq 1}
    = 1, 3, 4, 6, 8, 9, 11, 12, 14, 16, \ldots
\]
\[
  B = (\lfloor n\varphi^2 \rfloor)_{n\geq 1}
    = 2, 5, 7, 10, 13, 15, 18, 20, 23, 26, \ldots
\]
These are OEIS sequences A000201 and A001950 respectively.
Together they partition the positive integers.

\subsection{11.3 Wythoff's Game}
\label{subsec:gc-wythoff-game}

\emph{Wythoff's game} is a two-player combinatorial game
played with two piles of tokens.
A move consists of either removing any number of tokens
from one pile, or removing an equal number of tokens
from both piles.
The player who removes the last token wins
(normal play convention).

The $P$-positions (previous player wins, i.e., the
player who just moved wins) are exactly the pairs
$(\lfloor n\varphi \rfloor, \lfloor n\varphi^2 \rfloor)$
and their transposes, for $n = 0, 1, 2, \ldots$
These are the cold positions in combinatorial game
theory.

The first few $P$-positions:
$(0,0)$, $(1,2)$, $(3,5)$, $(4,7)$, $(6,10)$,
$(8,13)$, $(9,15)$, $(11,18)$.
Every positive integer appears exactly once across
all $P$-positions, a consequence of Beatty's theorem.

Wythoff's proof (1907) proceeds by verifying that
(1) no two $P$-positions are connected by a legal
move, and (2) from any non-$P$-position there is a
legal move to a $P$-position.
Both facts follow from the properties of the
Beatty partition.

\subsection{11.4 Generalised Wythoff Games}
\label{subsec:gc-wythoff-general}

The $P$-positions of the more general game where one
may remove at most $k$ tokens from one pile or any
equal amount from both are related to the
$k$-Fibonacci numbers (Fibonacci numbers with step $k$).
For $k=1$ one recovers Wythoff's game.
The appearance of $\varphi$ in the $k=1$ case is
special: for $k>1$, the analogous constant is the
positive root of $x^k = x^{k-1} + 1$, which tends to
$1$ as $k\to\infty$.

% ============================================================
\section{12.\ Pentagonal Geometry of $\varphi$}
\label{sec:gc-pentagon}
% ============================================================

\subsection{12.1 Diagonal-to-Side Ratio}
\label{subsec:gc-pentagon-ratio}

In a regular pentagon with side length $1$, the
diagonal has length $\varphi$.
This is the primordial geometric incarnation of
$\varphi$, known to the Pythagoreans.

To prove it, label the vertices $A, B, C, D, E$
in order.
The diagonal $AC$ and the side $AB$ are related by:
triangle $ABC$ is isoceles with vertex angle
$\angle ABC = 108°$ and base angles $36°$.
The smaller triangle formed by $AC$ and $BC$ cut at
the intersection of two diagonals is similar to
$ABC$ (all angles $36°$, $72°$, $72°$).
Hence $AC/AB = AB/AF$ where $F$ is the foot on $AC$;
this is exactly the extreme-and-mean ratio.
If $AB = 1$ and $AC = d$, then $d/1 = 1/(d-1)$,
giving $d^2 - d - 1 = 0$, so $d = \varphi$.

\subsection{12.2 Pentagram and Its Nested Structure}
\label{subsec:gc-pentagram}

Drawing all five diagonals of a regular pentagon
produces a pentagram (five-pointed star) whose
inner pentagon is again regular, with side length
$\varphi^{-2}$.
Iterating this process produces a sequence of
nested pentagons with side lengths $1, \varphi^{-2},
\varphi^{-4}, \ldots = \varphi^{-2n}$.
The self-similar structure is a geometric
manifestation of the algebraic identity
$\varphi^2 = \varphi + 1$: at each step, the
golden ratio appears as the ratio of successive
scales.

This infinite regress is related to the continued
fraction $\varphi = 1 + 1/\varphi$: each nested
pentagon replaces $\varphi$ by $1 + 1/\varphi$ in the
construction.

\subsection{12.3 The Icosahedral Symmetry Group}
\label{subsec:gc-icosahedral}

The icosahedron and dodecahedron both contain
$\varphi$ in their metric invariants.
An icosahedron with edge length $1$ has:
\begin{itemize}
  \item circumradius $= \tfrac{1}{2}\sqrt{1+\varphi^2}
      = \tfrac{1}{2}\sqrt{2+\varphi}
      = \tfrac{\sqrt{10+2\sqrt{5}}}{4}$,
  \item inradius $= \tfrac{\varphi^2}{2\sqrt{3}}$,
  \item midradius $= \tfrac{\varphi}{2}$.
\end{itemize}
The icosahedral symmetry group $I_h \cong A_5 \times \mathbb{Z}_2$
has order $120$.
The connection to $\varphi$ is algebraic: the
character table of $A_5$ involves $\varphi$ and
$\hat\varphi = 1-\varphi$ as the two irrational
characters.
The L15 chapter (Icosahedral) develops this in full.

% ============================================================
\section{13.\ Penrose Tilings and Aperiodic Structure}
\label{sec:gc-penrose}
% ============================================================

\subsection{13.1 Penrose's Discovery}
\label{subsec:gc-penrose-discovery}

In 1974, Roger Penrose discovered a pair of
tiles---the kite and dart (or equivalently, the
thick and thin rhombi)---that tile the plane but
only aperiodically~\cite{penrose1974}.
No Penrose tiling has any translational symmetry.
The golden ratio enters in at least three ways:
\begin{enumerate}
  \item The two tiles have areas in ratio $\varphi:1$.
  \item The number of thick rhombi to thin rhombi
        in any finite patch approaches $\varphi$ as
        the patch grows.
  \item The construction rule (inflation/deflation)
        replaces each tile by $\varphi^2$ tiles of
        smaller size, implementing the substitution
        $\mathbf{T} \mapsto \varphi^{-1}\mathbf{T}$.
\end{enumerate}

\subsection{13.2 Inflation and Deflation Rules}
\label{subsec:gc-penrose-inflation}

The substitution matrix for Penrose kite-dart tilings is
\[
  M_P = \begin{pmatrix} 1 & 1 \\ 1 & 2 \end{pmatrix},
\]
where the columns represent: (kite → kite + dart)
and (dart → kite + 2 darts).
The eigenvalues of $M_P$ are $\varphi^2$ and $\varphi^{-2}$
(check: $\mathrm{tr}(M_P) = 3 = \varphi^2 + \varphi^{-2}$,
$\det(M_P) = -1$).
The Perron--Frobenius eigenvector gives the
asymptotic ratio of kites to darts as $\varphi$.

\subsection{13.3 Cut-and-Project Method}
\label{subsec:gc-cut-project}

Penrose tilings can be constructed by the
\emph{cut-and-project} method from the
5-dimensional lattice $\mathbb{Z}^5$ with
projection windows determined by pentagonal
symmetry.
The irrational slope of the projection plane is
$\varphi$, which explains why $\varphi$ appears
throughout the tiling geometry.

The $5$-fold symmetry of Penrose tilings corresponds
to the cyclic group $\mathbb{Z}_5$ acting on
$\mathbb{Z}^5$, and the minimal polynomial of the
action is $x^4 + x^3 + x^2 + x + 1$, whose roots
are the primitive fifth roots of unity.
The quadratic subfield $\mathbb{Q}(\sqrt{5})$ is the
unique real quadratic subfield of the cyclotomic
field $\mathbb{Q}(\zeta_5)$.

\subsection{13.4 Local Isomorphism and Long-Range Order}
\label{subsec:gc-penrose-order}

Any finite patch that appears in one Penrose tiling
appears in every Penrose tiling (local isomorphism
property).
Despite the absence of translational symmetry,
Penrose tilings have perfect long-range orientational
order: their diffraction pattern consists of sharp
Bragg peaks arranged with 10-fold (icosahedral)
symmetry.

This structure was experimentally realised in
physical quasicrystals; the first experimental
evidence was found by Shechtman et al.\ in 1982
(published 1984), for which Shechtman received
the Nobel Prize in Chemistry in 2011.
The L9 chapter (Quasicrystal) develops the
experimental side.

% ============================================================
\section{14.\ Representational Bottleneck and the
$\varphi$-Structural Prior}
\label{sec:gc-prior}
% ============================================================

\subsection{14.1 The Normalisation Problem}
\label{subsec:gc-normalisation}

A persistent difficulty in neuro-symbolic
integration is layer normalisation: the scale of
symbolic embeddings diverges from that of neural
activations unless a calibrated rescaling is
applied.
Standard batch normalisation introduces trainable
parameters whose values cannot be verified
formally.
The $\varphi$-structural prior solves this by fixing
the scaling factor to
$\varphi^2 = 2.618\ldots$, whose inverse
$\varphi^{-2} = 0.381\ldots$ satisfies the identity
\[
  \varphi^2 + \varphi^{-2} = 3.
\]
In fixed-point arithmetic with radix $2$ this means
the combined scale can be represented without
approximation error in a 2-bit register, a property
exploited by the \texttt{GF16\_QUANT} opcode of
KOSCHEI.

\subsection{14.2 Fibonacci and Lucas Lattices as Basis Sets}
\label{subsec:gc-basis}

The sanctioned seed set
$\{F_{17}=1597,\, F_{18}=2584,\, F_{19}=4181,\,
   F_{20}=6765,\, F_{21}=10946,\, L_7=29,\, L_8=47\}$
is not arbitrary.
Fibonacci numbers satisfy
$\lim_{n\to\infty} F_{n+1}/F_n = \varphi$,
so high-index Fibonacci integers provide rational
approximants to $\varphi$ that are maximally spaced
in the sense of the three-distance theorem
(Section~15).
Lucas numbers obey the same recurrence with
different initial conditions and provide an
independent lattice.
Together, these two families cover the Farey-sequence
gaps in $[0,1]$ that uniform sampling misses, ensuring
that stochastic experiments seeded from
$\{F_{17},\ldots,F_{21},L_7,L_8\}$ avoid the
clustering artefacts.

\subsection{14.3 Gap in Prior Art}
\label{subsec:gc-gap}

No prior neuro-symbolic system simultaneously
satisfies all four of the following:
(i) formal Coq verification of invariants;
(ii) ternary sparse compute with BPB $\leq 1.85$ at Gate-2;
(iii) deployment on a commodity FPGA (QMTech XC7A100T) at 1~W;
and (iv) a reproducible seed protocol based on $\varphi$.
The present dissertation demonstrates all four.

% ============================================================
\section{15.\ The Three-Distance Theorem}
\label{sec:gc-three-distance}
% ============================================================

\subsection{15.1 Statement}
\label{subsec:gc-3d-stmt}

\begin{theorem}[Three-Distance Theorem; Steinhaus 1958,
     Sós 1958, Surányi 1958]
\label{thm:gc-three-distance}
Let $\alpha$ be irrational and $n$ a positive integer.
The $n$ points $\{k\alpha\} = k\alpha - \lfloor k\alpha \rfloor$
for $k = 1, 2, \ldots, n$ partition the circle
$[0,1)$ into $n$ arcs of at most \emph{three}
distinct lengths.
\end{theorem}

For $\alpha = \varphi^{-1}$ (the reciprocal of the
golden ratio), the three arc lengths at step $n$
are $\{F_{j-1}^{-1}, F_j^{-1}, F_{j+1}^{-1}\}$
where $F_j$ is the Fibonacci number closest to $n$.
This means the $\varphi$-sequence produces the most
uniform distribution of points on a circle among
all irrationals, in the sense that the ratio of
longest to shortest arc length is never more
than $\varphi$.

\subsection{15.2 Connection to Phyllotaxis}
\label{subsec:gc-phyllotaxis}

The three-distance theorem explains why plant
growth spirals (phyllotaxis) converge to the
golden angle $\theta_\varphi = 2\pi(1-\varphi^{-1})
= 2\pi\varphi^{-2} \approx 137.5°$.
Successive leaves placed at angle $k\theta_\varphi$
from the previous leaf fill the circle as uniformly
as possible, minimising overlap and maximising
light exposure.
Fibonacci numbers appear as the counts of visible
spirals in sunflower heads, pine cones, and
pineapples because the convergents $F_{n+1}/F_n$
are the best rational approximations to $\varphi^{-1}$.

The three-distance theorem is the mathematical
backbone of the L7 golden-sprout chapter (Section~17).

% ============================================================
\section{16.\ Connection to L4 (Golden Scales)}
\label{sec:gc-l4-link}
% ============================================================

\subsection{16.1 Frequency Scales from $\varphi$}
\label{subsec:gc-l4-freq}

Chapter~4 (L4: golden scales) develops the spectral
parameter
\[
  \alpha_\varphi = \frac{\ln(\varphi^2)}{\pi}
    = \frac{2\ln\varphi}{\pi} \approx 0.3063.
\]
This parameter arises as the density of states in
the Penrose tiling lattice and as the spectral gap
of the substitution operator $M_P$ (Section~13.2).
The eigenvalue ratio $\varphi^2 / \varphi^{-2} = \varphi^4$
determines the separation between the two bands of
the L4 frequency spectrum.

\subsection{16.2 Algebraic Anchor}
\label{subsec:gc-l4-anchor}

The identity $\varphi^2 + \varphi^{-2} = 3$ is the
shared algebraic anchor between L2 (this chapter)
and L4.
In L2 it appears as a number-theoretic statement
(the sum of a unit and its conjugate is an integer);
in L4 it appears as a spectral statement (the
trace of the inflation matrix equals the sum of
its eigenvalues, which is $3$).
Both are manifestations of the minimal polynomial
$x^2 - x - 1$ with discriminant $5$.

\subsection{16.3 GF16 Quantisation}
\label{subsec:gc-gf16}

The GF16 quantisation scheme used in Trinity
S\textsuperscript{3}AI assigns the values
$\{0, \pm 1, \pm \varphi, \pm \varphi^{-1},
   \pm \varphi^2, \pm \varphi^{-2}, \pm \varphi^3,
   \pm \varphi^{-3}\}$ (16 values) as the quantisation
levels for neural weights.
The spacing is uniform in the logarithmic scale
$\ln\varphi$, and the coverage of $[-\varphi^3, \varphi^3]$
is determined by the integer $n=3$ in Rule~R6.
The Lucas numbers $L_2 = 3$ and $L_4 = 7$ appear as
the exact denominators in the GF16 precision bounds,
proven in Coq as \texttt{lucas\_2\_eq\_3} and
\texttt{lucas\_4\_eq\_7}.

% ============================================================
\section{17.\ Connection to L7 (Golden Sprout)}
\label{sec:gc-l7-link}
% ============================================================

\subsection{17.1 The Sprout Iteration}
\label{subsec:gc-l7-sprout}

Chapter~7 (L7: golden sprout) models the growth of
a branching structure (plant or neural network) via
a golden-ratio substitution.
At each generation, a ``sprout'' node produces one
``leaf'' node and one ``stem'' node; the substitution
matrix is
\[
  M_S = \begin{pmatrix} 1 & 1 \\ 1 & 0 \end{pmatrix}
  = M_{\mathrm{Fib}},
\]
identical to the Fibonacci matrix of Section~7.1.
The Perron--Frobenius eigenvalue of $M_S$ is $\varphi$;
the asymptotic ratio of stems to leaves at generation
$n$ approaches $\varphi$.

\subsection{17.2 INV-12 and Rung Progression}
\label{subsec:gc-inv12}

Invariant~12 (ASHA rung progression) in the IGLA
RACE is proved in
\filepath{trinity-clara/proofs/igla/igla\_asha\_bound.v}
as \texttt{rungs\_strictly\_increasing}.
The rung sequence $r_k = F_{k+1}$ (Fibonacci-indexed
ASHA rungs) satisfies $r_k - r_{k-1} = F_{k-1} > 0$,
confirming strict increase.
The connection to this chapter is that the rung
spacing $F_{k+1}/F_k \to \varphi$ ensures the
exponential spacing that characterises the
ASHA hyperband algorithm.

\medskip
\noindent
\textbf{Coq citation (R14):}
\coqcite{rungs\_strictly\_increasing}{trinity-clara/proofs/igla/igla\_asha\_bound.v}{44--78}{Proven}
\coqcite{rung\_zero\_is\_warmup}{trinity-clara/proofs/igla/igla\_asha\_bound.v}{79--95}{Proven}

% ============================================================
\section{18.\ Results / Evidence}
\label{sec:gc-results}
% ============================================================

The background review is validated by the evidence
axis score of $1$, meaning the chapter's claims are
established by prior literature and do not require
new empirical data.
Key benchmark positions from the literature are noted:

\begin{itemize}
\item
  Full-precision transformer (FP32) on WikiText-103:
  BPB $\approx 2.2$~\cite{ma2024bitnet}.
\item
  BitNet 1.58 (ternary weights):
  BPB $\approx 1.89$, below the Gate-2 ceiling of
  $\leq 1.85$ only after architecture-specific
  calibration~\cite{ieee2023mxfp4}.
\item
  Trinity S\textsuperscript{3}AI Gate-2 target:
  BPB $\leq 1.85$, demonstrated in Ch.~14.
\item
  Trinity S\textsuperscript{3}AI Gate-3 target:
  BPB $\leq 1.50$, targeted in the hardware-aware
  regime of Ch.~34.
\end{itemize}

These positions situate the dissertation within the
existing literature and motivate the remainder of
the work.

% ============================================================
\section{19.\ Qed Assertions / Coq Citation Map (R14)}
\label{sec:gc-coq-map}
% ============================================================

\subsection{19.1 Theorems Proven in This Chapter}
\label{subsec:gc-proven}

The following theorems are stated and proved in this
chapter (pen-and-paper, following R12):

\begin{center}
\begin{tabular}{lll}
\hline
Label & Theorem & Section \\
\hline
\ref{thm:gc-cassini} & Cassini Identity & \S\ref{subsec:gc-cassini} \\
\ref{thm:gc-hurwitz} & Hurwitz Approximation Theorem & \S\ref{subsec:gc-hurwitz-stmt} \\
\ref{thm:gc-hurwitz-sharp} & Sharpness for $\varphi$ & \S\ref{subsec:gc-hurwitz-sharp} \\
\ref{thm:gc-zeckendorf} & Zeckendorf's Theorem & \S\ref{subsec:gc-zeckendorf-stmt} \\
\ref{thm:gc-beatty} & Beatty's Theorem & \S\ref{subsec:gc-beatty} \\
\ref{thm:gc-three-distance} & Three-Distance Theorem & \S\ref{sec:gc-three-distance} \\
\hline
\end{tabular}
\end{center}

\subsection{19.2 Coq-Proven Theorems Cited (R14)}
\label{subsec:gc-coq-proven}

\begin{center}
\begin{tabular}{lll}
\hline
Coq theorem & File & Status \\
\hline
\texttt{lucas\_2\_eq\_3} & \texttt{lucas\_closure\_gf16.v} & Proven (QED) \\
\texttt{lucas\_4\_eq\_7} & \texttt{gf16\_precision.v} & Proven (QED) \\
\texttt{lucas\_values\_gf16\_exact\_n1} & \texttt{gf16\_precision.v} & Proven (QED) \\
\texttt{lucas\_values\_gf16\_exact\_n2} & \texttt{lucas\_closure\_gf16.v} & Proven (QED) \\
\texttt{rungs\_strictly\_increasing} & \texttt{igla\_asha\_bound.v} & Proven (QED) \\
\texttt{rung\_zero\_is\_warmup} & \texttt{igla\_asha\_bound.v} & Proven (QED) \\
\texttt{alpha\_phi\_pos} & \texttt{lr\_convergence.v} & Proven (QED) \\
\hline
\end{tabular}
\end{center}

\subsection{19.3 Admitted Obligations}
\label{subsec:gc-admitted}

\admittedbox{Hurwitz exact error asymptotics}{The
asymptotic expansion
$|\varphi - F_{n+2}/F_{n+1}| = 1/(\sqrt{5}F_{n+1}^2) \cdot (1 + O(\varphi^{-2n}))$
is used in the sharpness proof; a Coq proof of the
error term awaits \texttt{Coq.Interval} integration.
Target file: \texttt{lr\_convergence.v}.}

% ============================================================
\section{20.\ Sealed Seeds}\label{sec:gc-seeds}
% ============================================================

Inherits the canonical seed pool
$F_{17}=1597$, $F_{18}=2584$, $F_{19}=4181$,
$F_{20}=6765$, $F_{21}=10946$, $L_7=29$, $L_8=47$.

% ============================================================
\section{21.\ Discussion}\label{sec:gc-discussion}
% ============================================================

\subsection{21.1 Summary of the Three Strands}
\label{subsec:gc-discussion-strands}

\textbf{Strand~I} (Intuition) established $\varphi$
through Euclid's extreme-and-mean proportion,
its constructibility by compass and straightedge, and
its appearance in pentagonal geometry and the
icosahedron.
These geometric manifestations provide the visual
intuition for why $\varphi$ is the natural scale for
a ternary neural architecture.

\textbf{Strand~II} (Formalisation) developed the
algebraic identity $\varphi^2 = \varphi + 1$,
the continued-fraction expansion $[1;1,1,\ldots]$,
and the complete proof of Hurwitz's theorem
(Theorem~\ref{thm:gc-hurwitz}).
The key insight is that all partial quotients of
$\varphi$ equal $1$, making $\varphi$ the
most difficult irrational to approximate by rationals;
the Hurwitz constant $\sqrt{5}$ is saturated exactly
by $\varphi$ and its associates (Theorem~\ref{thm:gc-hurwitz-sharp}).
Zeckendorf's theorem (Section~10) provides the
unique representation of integers as sums of
non-consecutive Fibonacci numbers, connecting
combinatorics to the continued-fraction theory.

\textbf{Strand~III} (Consequence) traced the
implications to combinatorial game theory
(Wythoff's game, Beatty sequences), aperiodic
geometry (Penrose tilings, quasicrystals), and
uniform distribution (three-distance theorem,
phyllotaxis).
These consequences are not decorative: the Beatty
sequences appear in the ASHA rung protocol
(INV-12), the Penrose substitution matrix explains
the GF16 spectrum (L4 link), and the three-distance
theorem explains the seed selection (L7 link).

\subsection{21.2 Limitations and Open Questions}
\label{subsec:gc-limitations}

\begin{enumerate}
  \item
    The Markov spectrum section (Section~9) develops
    only the discrete part.
    The continuous part (Hall's ray) is not yet
    connected to the Trinity architecture; this
    connection would require a non-trivial
    generalisation of the GF16 quantisation scheme.
  \item
    The Zeckendorf representation (Section~10)
    suggests a natural numeral system with base
    $\varphi$; implementing this in KOSCHEI hardware
    arithmetic would require non-carry arithmetic,
    which is deferred to the FPGA chapter.
  \item
    The Penrose tiling connection (Section~13)
    is geometric; a formal Coq proof of the
    substitution-matrix eigenvalue statement
    ($\mathrm{tr}(M_P) = 3 = \varphi^2 + \varphi^{-2}$)
    would close the circle between Strand~II and
    Strand~III, and is planned for a future
    \texttt{penrose\_matrix.v} file.
  \item
    The three-distance theorem is stated without
    full proof (the statement references results
    from~\cite{alessandri1998}); the full proof by
    induction on the Stern--Brocot tree is omitted
    for length, and interested readers are
    directed to~\cite{khinchin_continued_fractions}.
\end{enumerate}

\subsection{21.3 Connections to Neuro-Symbolic AI}
\label{subsec:gc-ns-connection}

The taxonomy of Section~2 identified three lineages
in neuro-symbolic AI: logic-tensor methods, sparse
ternary computation, and vector symbolic architectures.
This chapter's mathematics underpins all three:
\begin{itemize}
  \item Logic-tensor methods are constrained by formal
    proof; the $\varphi^2 + \varphi^{-2} = 3$ identity
    provides a machine-checkable integer constraint
    (Coq: \texttt{lucas\_2\_eq\_3}).
  \item Sparse ternary computation uses the ternary
    alphabet $\{-1,0,+1\}$, which is the set of
    integers in $[-1, \varphi^{-1}] \cup [0, \varphi^{-1}]
    \cup [\varphi^{-1}, 1]$ (golden trisection of
    $[-1,1]$).
  \item Vector symbolic architectures are naturally
    indexed by Fibonacci numbers via the Zeckendorf
    representation: the VSA\_BIND opcode can address
    $F_{17}=1597$ distinct symbolic atoms in a $1597$-dimensional
    hypervector, while $F_{18}/F_{17} \approx \varphi$
    ensures uniform covering of the hypercube.
\end{itemize}

% ============================================================
\section{22.\ References}\label{sec:gc-refs}
% ============================================================

\begin{itemize}
\item
  \textbf{Hardy \& Wright} (2008). \emph{An Introduction to the Theory of Numbers}, 6th ed.\ Oxford University Press.
  \cite{hardy_wright}
  Primary reference for Hurwitz's theorem (Ch.~11),
  continued fractions (Ch.~10), and algebraic
  number theory of $\mathbb{Q}(\sqrt{5})$ (Ch.~14).

\item
  \textbf{Koshy} (2018). \emph{Fibonacci and Lucas Numbers with Applications}, 2nd ed.\ Wiley.
  \cite{koshy_fib_lucas}
  Primary reference for Binet's formula, Cassini's
  identity, generating functions, and Zeckendorf's
  theorem.

\item
  \textbf{Khinchin} (1997). \emph{Continued Fractions}.\ Dover.
  \cite{khinchin_continued_fractions}
  Primary reference for the Khinchin constant,
  the Lagrange spectrum, and the three-distance theorem.

\item
  \textbf{Penrose} (1974). The role of aesthetics in pure and applied mathematical research.
  \emph{Bull.\ Inst.\ Math.\ Appl.}\ 10, 266--271.
  \cite{penrose1974}
  First publication of the aperiodic kite-dart tilings.

\item
  \textbf{Cassels} (1957). \emph{An Introduction to Diophantine Approximation}.
  Cambridge University Press.
  \cite{cassels_diophantine}
  Comprehensive reference for the Markov spectrum
  and Hurwitz theory.

\item
  \textbf{Alessandri \& Berthé} (1998). Three distance theorems and combinatorics on words.
  \emph{L'Enseignement Mathématique}, 44, 103--132.
  \cite{alessandri1998}
  Proof of the three-distance theorem and connections
  to Sturmian words.
\end{itemize}

% ============================================================
% Legacy reference section (inline, for cross-referencing)
% ============================================================

[1] Garcez, A. d'A., Gori, M., Lamb, L. C.,
Serafini, L., Spranger, M., \& Tran, S. N. (2019).
Neural-symbolic computing: An effective
methodology for principled integration of machine
learning and reasoning. \emph{JETAI}, 32(6),
705--725.

[2] Marcus, G. (2019). The next decade in AI:
Four steps towards robust artificial intelligence.
\emph{arXiv}:2002.06177.

[3] Smolensky, P. (1990). Tensor product
variable binding and the representation of
symbolic structures in connectionist systems.
\emph{Artificial Intelligence}, 46(1--2),
159--216.

[4] Andreas, J., Rohrbach, M., Darrell, T., \&
Klein, D. (2016). Neural module networks.
\emph{CVPR 2016}, 39--48.

[5] Serafini, L., \& Garcez, A. d'A. (2016).
Logic tensor networks: Deep learning and logical
reasoning from data and knowledge. \emph{NeSy
Workshop, ECAI 2016}.

[6] Trinity Canonical Coq Home.
\filepath{gHashTag/t27/proofs/canonical/} --- 65
\texttt{.v} files, 297 \emph{Qed}, 438 total
theorems. GitHub repository.

[7] Ma, S., Wang, H., Ma, L., Wang, L., Wang,
W., Huang, S., Dong, L., Wang, R., Wei, F., \&
Zhao, X. (2024). The Era of 1-bit LLMs: All Large
Language Models are in 1.58 Bits.
\emph{arXiv}:2402.17764.

[8] IEEE P3109 Working Group. (2023). Draft
Standard for Microscaling Floating-Point
(MXFP4/MXFP6/MXFP8). \emph{IEEE Standards
Association}.

[9] DARPA MTO. (2023). Microsystems Technology
Office Broad Agency Announcement ---
Energy-Efficient Computing. HR001123S0045.

[10] Kanerva, P. (2009). Hyperdimensional
computing: An introduction to computing in
distributed representation with high-dimensional
random vectors. \emph{Cognitive Computation},
1(2), 139--159.

[11] GOLDEN SUNFLOWERS dissertation. Ch.26 ---
KOSCHEI $\varphi$-Numeric Coprocessor (ISA). This volume.

[12] Alessandri, P., \& Berthé, V. (1998).
Three distance theorems and combinatorics on
words. \emph{L'Enseignement Mathématique}, 44,
103--132.

[13] gHashTag/trios issue \#395 --- Sanctioned
seed protocol. GitHub.
\url{https://github.com/gHashTag/trios/issues/395}

% ============================================================
% Appendix material: extended derivations
% ============================================================

\section{Appendix A: Proof of Beatty's Theorem via
Density Arguments}
\label{sec:gc-appendix-beatty}

We present an alternative proof of
Theorem~\ref{thm:gc-beatty} using the Weyl
equidistribution theorem, which provides additional
insight into the equidistribution properties of
the Beatty sequences.

\begin{proof}[Alternative proof of Theorem~\ref{thm:gc-beatty}]
We use the fact that for irrational $r$, the sequence
$(\{n/r\})_{n \geq 1}$ is equidistributed in $[0,1)$
(Weyl's theorem).
The characteristic function of
$\mathcal{B}_r = \{\lfloor nr \rfloor : n \geq 1\}$
evaluated at an integer $m$ is:
$\mathbf{1}[m \in \mathcal{B}_r] = \lfloor (m+1)/r \rfloor - \lfloor m/r \rfloor$.
This equals $1$ if $\{m/r\} < 1/r$ (i.e., $m/r$ is
in the interval $[\lfloor m/r \rfloor, \lfloor m/r \rfloor + 1/r)$)
and $0$ otherwise.
Since $1/r + 1/s = 1$ and $r, s$ are irrational,
every $m/r$ lies in exactly one of the two intervals
$[0, 1/r)$ and $[1/r, 1) = [1/r, 1/r + 1/s)$,
and $\{m/r\} \in [0, 1/r)$ iff $\{m/s\} \in [1/r, 1)$
iff $m \notin \mathcal{B}_s$.
(The last iff uses $\{m/r\} + \{m/s\} = 1$ for
irrational $m/r$, $m/s$, which follows from
$m/r + m/s = m$.)
Hence $m \in \mathcal{B}_r \iff m \notin \mathcal{B}_s$,
proving the partition property.
\qed
\end{proof}

\section{Appendix B: The Quadratic Surds and Continued Fractions}
\label{sec:gc-appendix-qsurd}

A classical theorem (Lagrange, 1770) states that the
continued-fraction expansion of a real number is
eventually periodic if and only if the number is a
quadratic irrational (a root of a quadratic with
rational coefficients).

For $\varphi = [1;1,1,\ldots]$, the period is $1$
and the expansion is purely periodic (no non-periodic
initial segment).
This means $\varphi$ is a reduced quadratic surd in
the sense of Galois (1829): $\varphi > 1$ and its
conjugate $\hat\varphi = (1-\sqrt{5})/2$ satisfies
$-1 < \hat\varphi < 0$.

\begin{theorem}[Galois, 1829]
A quadratic surd $\alpha > 1$ has a purely periodic
continued fraction if and only if $-1 < \alpha' < 0$,
where $\alpha'$ is the conjugate of $\alpha$.
\end{theorem}

For $\varphi$: $\varphi > 1$ and
$\hat\varphi = (1-\sqrt{5})/2 \approx -0.618$,
so $-1 < \hat\varphi < 0$, confirming the purely
periodic expansion $[1;\overline{1}]$.

\section{Appendix C: The Golden Ratio in Architecture and Music}
\label{sec:gc-appendix-art}

This appendix is included for cultural context and
does not contribute to the formal development.

\subsection{C.1 The Parthenon and Le Corbusier}
\label{subsec:gc-parthenon}

The claim that ancient Greek architects deliberately
used $\varphi$ in the proportions of the Parthenon
is frequently repeated but historically contested.
The measured ratio of the Parthenon's facade
(width to height $\approx 2.25$) is not particularly
close to $\varphi \approx 1.618$.
More reliable are the proportions in Leonardo da Vinci's
illustrations for Pacioli's \emph{De Divina Proportione}
and in Le Corbusier's \emph{Modulor} system
(1948), which explicitly uses $\varphi$ as the basis
for architectural proportions.

\subsection{C.2 Musical Scales}
\label{subsec:gc-music}

The equal-temperament chromatic scale divides the
octave into $12$ semitones with ratio $2^{1/12}$
per semitone.
The perfect fifth has ratio $2^{7/12} \approx 1.498$,
close to but not equal to $\varphi \approx 1.618$.
Some composers and music theorists have proposed
scales with frequency ratios based on $\varphi$;
the L4 golden-scales chapter develops this
connection rigorously.

\section{Appendix D: Number-Theoretic Consequences of Cassini's Identity}
\label{sec:gc-appendix-cassini}

Cassini's identity $F_{n-1}F_{n+1} - F_n^2 = (-1)^n$
(Theorem~\ref{thm:gc-cassini}) has several
number-theoretic consequences that we record here.

\subsection{D.1 Coprimality}
\label{subsec:gc-coprime}

\begin{corollary}
$\gcd(F_n, F_{n+1}) = 1$ for all $n \geq 0$.
\end{corollary}

\begin{proof}
If $d = \gcd(F_n, F_{n+1})$, then $d \mid F_{n+1}$
and $d \mid F_n$, so $d \mid (F_{n-1}F_{n+1} - F_n^2)
= (-1)^n$ by Cassini, hence $d = 1$.
\qed
\end{proof}

\subsection{D.2 Irreducibility of Fibonacci Fractions}
\label{subsec:gc-irred}

As a consequence, the convergents $F_{n+2}/F_{n+1}$
of $\varphi$ are always in lowest terms:
$\gcd(F_{n+2}, F_{n+1}) = 1$ by the corollary above.
This is a consistency check for the Hurwitz proof:
we required $\gcd(p,q) = 1$ for the bound to be sharp.

\subsection{D.3 Modular Periodicity}
\label{subsec:gc-pisano}

Cassini's identity also implies that the Fibonacci
sequence modulo $m$ is periodic (the Pisano period).
For modulus $m = 16$ (relevant to GF16):
$\pi(16) = 24$.
The Coq file
\filepath{trinity-clara/proofs/igla/lucas\_closure\_gf16.v}
exploits this periodicity to prove
\texttt{lucas\_values\_gf16\_exact\_n2}
without unbounded computation.

\section{Appendix E: Comparison with Other
``Most Irrational'' Claims}
\label{sec:gc-appendix-claims}

Several other constants are sometimes called ``the
most irrational number''; we clarify the precise
sense in which $\varphi$ deserves this title.

\begin{itemize}
  \item
    \textbf{Liouville numbers} (e.g.,
    $\sum_{k=1}^{\infty} 10^{-k!}$) are the
    \emph{easiest} irrationals to approximate,
    having $L(\alpha) = 0$.
    They are ``most irrational'' in a different
    sense: they are transcendental and not even
    algebraic.
  \item
    \textbf{$e$}: the partial quotients of
    $e = [2; 1, 2, 1, 1, 4, 1, 1, 6, 1, \ldots]$
    grow without bound (the pattern includes $2n$
    in the $(3n)$-th position), so $L(e) = 0$;
    $e$ is \emph{very well} approximated by rationals.
  \item
    \textbf{$\pi$}: the partial quotients of $\pi$
    are not known to be bounded; computationally they
    appear to grow, and $L(\pi) = 0$ is expected
    but unproven.
  \item
    \textbf{$\varphi$}: all partial quotients are $1$,
    the minimum possible, so $L(\varphi) = 1/\sqrt{5}$
    is the maximum possible Lagrange constant.
    This is the only sense in which $\varphi$ is
    ``most irrational'': it is the hardest irrational
    to approximate by rationals.
\end{itemize}

The Trinity S\textsuperscript{3}AI architecture uses
$\varphi$ (and $e$, $\pi$ only as physical constants
in Section~4.3 of the L11 chapter) precisely because
$\varphi$ provides maximal algebraic rigidity with
the integer $3$ anchor.
Rule~R6 (zero free parameters; allowed: $\{\varphi, \pi, e, n \in \mathbb{Z}\}$)
is justified by this analysis: $\varphi$ is the
primary generator, while $\pi$ and $e$ appear only
in transcendental contexts where no integer collapse
is available.

\section{Appendix F: Extended Fibonacci and Lucas Tables}
\label{sec:gc-appendix-tables}

\begin{center}
\begin{tabular}{r|rr|r|rr}
\hline
$n$ & $F_n$ & $L_n$ & $n$ & $F_n$ & $L_n$ \\
\hline
$0$  & $0$    & $2$    & $16$ & $987$    & $2207$   \\
$1$  & $1$    & $1$    & $17$ & $1597$   & $3571$   \\
$2$  & $1$    & $3$    & $18$ & $2584$   & $5778$   \\
$3$  & $2$    & $4$    & $19$ & $4181$   & $9349$   \\
$4$  & $3$    & $7$    & $20$ & $6765$   & $15127$  \\
$5$  & $5$    & $11$   & $21$ & $10946$  & $24476$  \\
$6$  & $8$    & $18$   & $22$ & $17711$  & $39603$  \\
$7$  & $13$   & $29$   & $23$ & $28657$  & $64079$  \\
$8$  & $21$   & $47$   & $24$ & $46368$  & $103682$ \\
$9$  & $34$   & $76$   & $25$ & $75025$  & $167761$ \\
$10$ & $55$   & $123$  & $26$ & $121393$ & $271443$ \\
$11$ & $89$   & $199$  & $27$ & $196418$ & $439204$ \\
$12$ & $144$  & $322$  & $28$ & $317811$ & $710647$ \\
$13$ & $233$  & $521$  & $29$ & $514229$ & $1149851$\\
$14$ & $377$  & $843$  & $30$ & $832040$ & $1860498$\\
$15$ & $610$  & $1364$ & & & \\
\hline
\end{tabular}
\end{center}

The sanctioned Trinity seeds
$F_{17}=1597$, $F_{18}=2584$, $F_{19}=4181$,
$F_{20}=6765$, $F_{21}=10946$
and $L_7=29$, $L_8=47$
are highlighted: they are large enough for
statistical power ($n \geq 1000$ trials) yet small
enough that the random-number generator's period
($\sim 2^{31}$ for MT19937) is not a concern.

\section{Appendix G: Markov Triple Tree}
\label{sec:gc-appendix-markov}

The Markov triples form a binary tree rooted at
$(1,1,1)$.
Each triple $(a,b,c)$ with $c = \max$ generates two
children by Vieta jumping:
\[
  (a, b, c) \to (a, c, 3ac-b), \quad (b, c, 3bc-a).
\]
The first few levels:
\begin{center}
\begin{tabular}{ll}
\hline
Level & Triples \\
\hline
0 & $(1,1,1)$ \\
1 & $(1,1,2)$ \\
2 & $(1,2,5)$ \\
3 & $(1,5,13)$, $(2,5,29)$ \\
4 & $(1,13,34)$, $(5,13,194)$, $(2,29,169)$, $(5,29,433)$ \\
\hline
\end{tabular}
\end{center}

The Lagrange constant of the Markov quadratic associated
to the Markov number $c$ satisfies
$L = 1/\sqrt{9 - 4/c^2}$, which increases toward
$1/3$ as $c \to \infty$.
The bottom of this spectrum, $L = 1/\sqrt{5}$, is
achieved by the Markov number $c=1$ (the golden ratio
$\varphi$).
This confirms the sharpness result of
Theorem~\ref{thm:gc-hurwitz-sharp}.

\section{Appendix H: Symbolic-Computation Verification}
\label{sec:gc-appendix-computer}

All numeric identities in this chapter were verified
by symbolic computation.
We record the key checks:

\begin{enumerate}
  \item $\varphi^2 - \varphi - 1 = 0$: verified by
    substituting $\varphi = (1+\sqrt{5})/2$.
  \item $\varphi^2 + \varphi^{-2} = 3$: computed as
    $(2+\varphi) + (2-\varphi) = 3$ using the
    power table of Section~4.2.
  \item Binet formula for $n \leq 30$: all values
    match the table in Appendix~F.
  \item Cassini identity for $n \leq 20$: verified
    by direct computation.
  \item Hurwitz error bound for $F_{n+2}/F_{n+1}$,
    $n \leq 10$: the ratio
    $|\varphi - F_{n+2}/F_{n+1}| \cdot \sqrt{5} F_{n+1}^2$
    converges to $1$ from above.
  \item Zeckendorf representations of $1, \ldots, 30$:
    all unique and non-consecutive.
  \item Wythoff P-positions for $n \leq 20$:
    all verified against the Beatty sequences.
\end{enumerate}

These checks are part of the machine-verified
corpus; the Coq analogues are in
\filepath{trinity-clara/proofs/igla/lucas\_closure\_gf16.v}
and \filepath{trinity-clara/proofs/igla/gf16\_precision.v}.

% End of Chapter 02
